Large language model (LLM) decoding contains a mix of compute-sensitive and memory-sensitive regions, which makes processing-in-memory (PIM) style offloading potentially attractive. The practical difficulty is that offload value is highly condition-dependent: some decode runs benefit substantially, while others do not. We present a trace-driven framework for studying \emph{virtual} PIM decision signals during LLM decoding. Our approach collects runtime-grounded decode traces from open-source stacks, attaches a prior-work-grounded virtual PIM cost model, and evaluates static and dynamic scheduling rules without claiming real hardware speedups. Using CPU and GPU traces across older GPT-style models and newer 4-bit instruct models, we find a consistent weak-signal versus positive-signal split: short-context decoding often shows no realistic gain, while longer-context decoding repeatedly becomes favorable for dynamic offload. Dense context sweeps show that context length is the clearest coarse regime-entry variable in our current setup, but the onset differs across model families. We further show that explicit KV-cache pressure proxies are more consistent with the observed boundary than a static attention heuristic alone, and that KV-aware features can help recover some early-positive families that a generic realistic feature policy misses. These results position virtual PIM not as a uniformly beneficial accelerator, but as a regime-dependent opportunity whose value can be characterized from runtime traces.
